\cotefile{1}

\scanpage{125}
\[
=\bigl(\alpha_n(|A|+|B|)\bigr)^{\unsure{1/2}}\bigl(\alpha_n(|A|'+|B|')^{1/2}
\]
\[
\deleted{\det(1+A+B)\leq\det(1+|A|)\det(1+|B|)}
\]

\fbox{Corollaire 3}
\[
|\det(1+A+B)|\leq\tfrac12\bigl(\det(1+|A|+|B|)+\det(1+|A^{\unsure{*}}|^{\unsure{\prime|,}}+|B^{\unsure{*}}|')
\]
résulte aussitôt de
\[
|\alpha_n(A+B)|\leq\tfrac12\bigl[\alpha_n(|A|+|B|)+\alpha_n(|A|'+|B|')\bigr]
\]

\fbox{Corollaire 4}
\[
\boxed{|\det(1+A+B)|\leq\det(1+|A|)\det(1+|B|)}
\]
Il suffit de \deleted{\unsure{prouver}} le prouver pour $A$ et $B$
hermitiens positifs\,; car alors on aura, d'après
le corollaire 3,
\[
\begin{aligned}
|\det(1+A+B)|\leq\tfrac12\bigl[&\det(1+|A|)\det(1+|B|)\\
&+\det(1+|A|')\det(1+|B|')\bigr],
\end{aligned}
\]
ce qui donné l'inégalité voulue, en vertu de
$\det(1+|A|)=\det(1+\unsure{\lvert A(\,{}'|\lvert A\rvert'}),\ \det(1+|B|)=\det(1+|B|')$

-- Quant au cas hermitien \inserted{positif}, on écrit
\[
\det(1+\unsure{R|h}+k)=\det(1+h)\bigl(1+\tfrac{h}{1+h}\bigr)=\det(1+h)\det(1+Hk)
\]
où $H=\frac{1}{1+h}$\,; on est ramené à prouver
$\det(1+Hk)\leq\det(1+k)$, il suffit pour cela
de prouver $\alpha_n(Hk)\leq\alpha_n(k)$ pour tout $n$,
soit $\operatorname{Tr}\overset{n}{\bigwedge}Hk\leq\|\overset{n}{\bigwedge}k\|_1$. Or on a
\[
\overset{n}{\bigwedge}Hk=(\overset{n}{\bigwedge}H)(\overset{n}{\bigwedge}k),\quad\text{d'où}
\]
\[
\|\overset{n}{\bigwedge}Hk\|_1\leq\|\overset{n}{\bigwedge}H\|_\infty\|\overset{n}{\bigwedge}k\|_1
\quad\text{et}\quad\|\overset{n}{\bigwedge}H\|_\infty\leq\|H\|_\infty^{\unsure{a|n}}
\]
\margin{!}or $\|H\|_\infty\leq1$.

\[
\boxed{\text{\fbox{Corollaire 5}}\quad\det(1+|A+B|)\leq\det(1+|A|)\det(1+|B|)}
\]

\scanpage{126}
\[
=\bigl(\alpha_n(|A|+|B|)\bigr)^{\unsure{1/2}}\bigl(\alpha_n(|A|'+|B|')^{1/2}
\]
\[
\deleted{\det(1+A+B)\leq\det(1+|A|)\det(1+|B|)}
\]

\fbox{Corollaire 3}
\[
|\det(1+A+B)|\leq\tfrac12\bigl(\det(1+|A|+|B|)+\det(1+|A^*|^{\unsure{\prime|,}}+|B^*|')\bigr.
\]
résulte aussitôt de
\[
|\alpha_n(A+B)|\leq\tfrac12\bigl[\alpha_n\unsure{(|A|+|B|)}+\alpha_n\unsure{(|A|'+|B|')}\bigr]
\]

\fbox{Corollaire 4}
\[
\boxed{|\det(1+A+B)|\leq\det(1+|A|)\det(1+|B|)}
\]
Il suffit de \deleted{\unsure{prouver}} le prouver pour $A$ et $B$
hermitiens positifs\,; car alors on aura, d'après
le corollaire 3,
\[
\begin{aligned}
|\det(1+A+B)|\leq\tfrac12\bigl[&\det(1+|A|)\det(1+|B|)\\
&+\det(1+|A|')\det(1+|B|')\bigr],
\end{aligned}
\]
ce qui donné l'inégalité voulue, en vertu de
$\det(1+|A|)=\det(1+\unsure{\lvert A(\,{}'|\lvert A\rvert'}),\ \det(1+|B|)=\det(1+|B|')$

-- Quant au cas hermitien \inserted{positif}, on écrit
\[
\det(1+\unsure{R|h}+k)=\det(1+h)\bigl(1+\tfrac{h}{1+h}\bigr)=\det(1+h)\det(1+Hk)
\]
où $H=\frac{1}{1+h}$\,; on est ramené à prouver
$\det(1+Hk)\leq\det(1+k)$, il suffit pour cela

Ce corollaire contient le précédent, car on sait
que $\det(1+\deleted{\unsure{|}}A+B)\leq\det(1+|A+B|)$. -- Il
résulte du précédent, car on aura ``en général''
$|A+B|=U(A+B)$, $U$ unitaire, d'autre part
\[
\begin{aligned}
\det(1+UA+UB)&\leq\det(1+|UA|)\det(1+|UB|),\\
&=\det(1+|A|)\det(1+|B|)\qquad(\text{car } |UA|\\
&=U|A|U^{-1},\quad |UB|=U|B|U^{-1})
\end{aligned}
\]

\scanpage{127}
\uline{Théorème} \unsure{Le déterminant} de Fredholm de $\unsure{u(\in L^1)|u\in L^1}$ \illegible{1}

\unsure{Ntn|Notn} $\mathcal L_0$\,; dire \inserted{\unsure{est-ce que}} \unsure{exactement} quelles
\unsure{parties|partie} du \deleted{\unsure{th}} théorème \unsure{--} sont connues.

\uline{Théorème}

1) Pour \deleted{\unsure{tout}} $0<p\leq\infty$, $\mathcal L^p(\mathcal H)$ est un idéal \inserted{bilatère}
dans $L(\mathcal H)$, stable pour l'involution $A\to A^*$. De
\deleted{\unsure{On a}} façon précise :
\[
\deleted{S_p}\|A\|_p=\deleted{S_p}\|A^*\|_p \tag*{(1)}
\]
\[
\begin{aligned}
\deleted{S_p}\|BA\|_p&\leq\|B\|\deleted{S_p}\|A\|_p
\qquad\text{pour } A\in\mathcal L^p(\mathcal H),B\in L(\mathcal H)\\
\|AB\|_p&\leq\|B\|\ \|A\|_p
\end{aligned}\tag*{(2)}
\]
Si $\infty\geq p\geq1$, $A\to\|A\|_p$ est une norme sur $\mathcal L^p(\mathcal H)$ \deleted{\unsure{(2)|(top.}}
i.e. on a
\[
\|A+B\|_p\leq\|A\|_p+\|B\|_p\quad\|\underset{\text{(1)}}{\lambda} A\|_p=|\lambda|\|A\|_p
\quad(1\leq p\leq\infty\tag*{(3)}
\]
\inserted{et $\mathcal L^p(\mathcal H)$ est complet pour cette norme}
Si $0<p\leq\infty$, alors l'expression $S_p(A)=\|A\|_p^p$ sur
$\mathcal L^p(A)$ satisfait à l'inégalité du triangle
\[
S_p(A+B)\leq S_p(A)+S_p(B)\qquad(0<p\leq1)\tag*{(4)}
\]
et muni de la distance, invariante par \unsure{translat.|transl.}
$S_p(A-B)$, $\mathcal L^p(A)$ est un espace vectoriel \unsure{top. log.|topolog.}
métrisable et complet, non localement convexe
si $p<1$, et $\dim\mathcal H=+\infty$.
\margin{\textcircled{2} \uline{Dual de $\mathcal L^p$}\,; \uline{réflexivité}}

\deleted{2}3) Si $A\in\mathcal L^p(\mathcal H)$ $(0<p<+\infty)$, alors la suite
des valeurs propres \inserted{$\alpha_i$} de $A$ est de puissance
$p^{\mathrm{ième}}$ sommable, et on a
\[
\sum|\alpha_i|^p\leq S_p(A)\tag*{(5)}
\]
\margin{généralisation aux normes de Schatten? \unsure{Ici plutôt|Ceci plutôt} en dernier \uline{Corollaire}}
$S_p(A)$ et $\|A\|_p$ sont des fonctions croissantes
sur l'ensemble des opérateurs hermitiens
positifs, i.e $H$ et $K$ hermitiens positifs, $H\leq K$,
implique $S_p(H)\leq S_p(K)$.

\deleted{\unsure{3|2}}\unsure{4|3}) Soient $0<p,q,r\leq+\infty$, avec $\frac{1}{p}+\frac{1}{q}\geq\frac{1}{r}$. Si \deleted{\unsure{donc}}
$A\in\mathcal L^p(\mathcal H)$, $B\in\mathcal L^q(\mathcal H)$, alors $AB\in\mathcal L^r(\mathcal H)$, et on a
\[
\|AB\|_r\leq\|A\|_p\|B\|_q\tag*{(6)}
\]
\margin{(1) .....}

\scanpage{128}
\margin{\fbox{\parbox{1.9cm}{Montrer qu'on \unsure{peut se ramener} à la dimension finie}}}
La première partie du théorème est résumée dans les
formules (1) à (4), \deleted{et} \inserted{outre} les affirmations relatives
aux topologies vectorielles des espaces $\mathcal L^p(\mathcal H)$.
\inserted{Montrons \deleted{\unsure{d'abord}} d'abord que ces dernières sont conséquences très faciles \deleted{\unsure{de|des}} \unsure{des formules} (3) et (4)\ldots}
La formule (1) est évidente, \unsure{car|La} \deleted{\unsure{on sait que}}
\deleted{\unsure{bien} \unsure{au} cas $p<+\infty$, et} \deleted{\unsure{pour $p<+\infty$}}, on sait (théorème
de Riesz) que $|A|$ et $|A^*|$ sont \deleted{\unsure{de même suite}} \inserted{unitairement équivalents}
(?) \deleted{\unsure{valeurs propres}}. La formule (2) est un cas
particulier de la formule (6), mais \deleted{\unsure{peut se|peut être}} \inserted{\unsure{on peut|Je peux}}
donner une démonstration \deleted{\unsure{pourtant directe}}
\deleted{démontrer simplement de façon directe}
directe, en procédant comme dans $[\,]$, .

En effet, si $(\unsure{k_i|\alpha_i})$ \deleted{\unsure{resp\textsuperscript{t} $(\beta_i)$}} désigne la
suite des valeurs \inserted{propres} de \deleted{$\unsure{\lvert BA\rvert|\lvert BA\rvert^2}$} $|A|$
\inserted{et $(\unsure{\beta_i|h_i})$ la suite des valeurs propres de $\unsure{\lvert BA\rvert|\lvert B\rvert A\rvert|\lvert BA\rvert^2}$ (resp\textsuperscript{t} $|AB|$)} \deleted{\unsure{resp\textsuperscript{t} $|BA|$}},
rangées par ordre de grandeur
décroissante (chacune répétée autant de fois
que le comporte sa multiplicité) on a
\[
\beta_i\leq\|B\|\alpha_i\quad\text{pour tout }i.\quad(\text{loc cité, page \quad, lemme \quad}),
\]
d'où aussitôt les formules (2).

La \deleted{\unsure{dernière} des formules (3) est évidente, la}
Les formules (3) \deleted{peut} \inserted{peuvent} être regardées \deleted{de la formule} comme
\uline{première} \deleted{\unsure{est une conséquence immédiate}}
\unsure{conséquences|conséquence} de la formule (6), (que nous démontrerons de façon indépendante) de la façon
suivante. Si $p'$ est tel que $\frac{1}{p}+\frac{1}{p'}=1$,
et si $B\in\mathcal L^{p'}(\mathcal H)$, \deleted{et tel que}
\inserted{$\|B\|_{p'}\leq1$}, alors
on a en vertu de (6) $\|AB\|_1\leq\|A\|_p$, d'où
\[
|\operatorname{Tr}(AB)|\leq\|A\|_p,\quad\text{et}\quad
\operatorname*{Sup}_{\|B\|_{p'}\leq1}|\operatorname{Tr}AB|\leq\|A\|_p
\]

\scanpage{129}
Cette inégalité peut être remplacée par une égalité,
car il existe un \unsure{op\textsuperscript{r}|op.} part\textsuperscript{t} is. $U$ tel que
$UA=|A|$, \deleted{\unsure{puis, une|puis, comme}} et d'autre part, comme il
résulte aussitôt de la décomposition spectral de
$|A|$ et du th. de Hölder classique sur les
espaces $L^p$ et $L^{p'}$, un opérateur hermitien
positif $H\in L^{p'}$ tel que $\operatorname{Tr}(H|A|)=\|A\|_p$,
soit \deleted{\unsure{$|$}}$\operatorname{Tr}(HUA)$\deleted{\unsure{$|$}}$=\|A\|_p$. Si on pose
$B=UA$, on aura $\|B\|_{p'}\leq\|U\|\ \|H\|_{p'}$ (formule
(1)) donc $\|B\|_{p'}\leq1$, ce qui prouve notre assertion.
Par suite on a bien
\[
\|A\|_p=\operatorname*{Sup}_{\|B\|_{p'}\leq1}|\operatorname{Tr}AB|\tag*{(7)}
\]
ce qui prouve que $\|A\|_p$ est bien une
semi norme, comme borne supérieure de
semi-normes $A\to|\operatorname{Tr}AB|$

Une telle démonstration n'est plus pratiquable
pour $p<1$\,; \deleted{et} nous démontrerons la formule (4)
par voie analytique, grace au

\uline{Lemme} 1 \quad Soit $(\alpha_i)$ une suite \deleted{décroissante} \inserted{\deleted{\unsure{de puiss}}} de
\margin{\rule[-3.9cm]{0.4pt}{4.1cm}}nombres $\geq0$, \deleted{\unsure{de puissance $p$\textsuperscript{ième} sommables}}
\inserted{et soit} \deleted{\unsure{Soit|soit}} $0<p<1$. Posons $f(\unsure{\bar z|z})=\prod_{i=1}^{\infty}(1+\alpha_i z)$.
Alors on a
\[
\sum\alpha_i^p=p\frac{\sin\pi p}{\pi}\int_0^\infty\frac{\log f(\unsure{r|v})}{r^{1+p}}\,dr\tag*{(8)}
\]
(les deux membres pouvant éventuellement
être infinis)

\scanpage{130}
\uline{Corollaire}\quad\margin{\rule[-3.8cm]{0.4pt}{3.8cm}} Si $A\in\mathcal L^1(\mathcal H)$, soit
$M_A(\unsure{v|r})=\operatorname*{Sup}_{|z|=\unsure{v|r}}\det(1+zA)$
\deleted{si $A$ est \unsure{hermitien} \unsure{positif}, et si $A\in\mathcal L^p(\mathcal H)$}, \deleted{\unsure{tel|on}}
\deleted{\unsure{avec}} \inserted{et} $0<p<1$\,; on a
\[
S_p(A)=p\frac{\sin\pi p}{\pi}\int_0^\infty\frac{\log M_{|A|}(\unsure{v|r})}{{\unsure{v|r}}^{1+p}}\,d\unsure{v|r}\tag*{(9)}
\]

Il suffit en effet de poser $f(z)=\det(1+z|A|)$,
les $\alpha_i$ étant les valeurs propres de $|A|$\,; on aura
ici $M(\unsure{v|r})=f(\unsure{v|r})$.

Pour démontrer (4), on \deleted{\unsure{pe|peut}} peut écrire
\[
S_p(A+B)=S_p(|A+B|)=p\frac{\sin\pi p}{\pi}\int_0^\infty
\frac{\log M_{|A+B|}(\unsure{v|r})}{{\unsure{v|r}}^{1+p}}\,d\unsure{v|r}
\]
Or $M_{|A+B|}(\unsure{v|r})=\det(1+\unsure{v|r}|A+B|)\leq\det(1+\unsure{v|r}|A|)\det(1+\unsure{v|r}|B|)$
(formule \unsure{\ldots}) \deleted{$\leq\unsure{rM_A|M_A}$} soit
\[
M_{|A+B|}(\unsure{v|r})\leq M_{|A|}(\unsure{v|r})\cdot M_{|B|}(\unsure{v|r})\tag*{(10)}
\]
d'où
\[
S_p(A+B)\leq\frac{p\sin\pi p}{\pi}\int_0^\infty\frac{\log M_{|A|}(\unsure{v|r})}{{\unsure{v|r}}^{1+\unsure{v|r|p}}}\,d\unsure{v|r}
+\frac{p\sin\pi p}{\pi}\int_0^\infty\frac{\log M_{|B|}(\unsure{v|r})}{{\unsure{v|r}}^{1+\unsure{r|v|p}}}\,d\unsure{v|r}
\]
Appliquant de nouveau (9), on obtient bien la
formule voulue (4).

\deleted{\unsure{Pour prouver}}

2) Pour prouver la formule (5), je ne connais
qu'une méthode analytique. De façon générale,
\unsure{.} si $A$ est un opérateur cp\textsuperscript{t} continu, on prend la
suite des val. propres $(\alpha_i)$, \deleted{\unsure{prises|posons}} et si
$0<p<+\infty$, posons $s_p(A)=\sum|\alpha_i|^p$.
\deleted{\unsure{On a} $\unsure{s_p(A)|S_p(A)}\unsure{\leq|=}$ \unsure{$+\infty$|$S_p$} \unsure{évidemment pour $p=1$|évidemment pour tout entier} Si $n>0$ $s_p(A)=s_{\frac{p}{n}}(A^n)$. \unsure{Supposons}}

\scanpage{131}
Supposons d'abord $p<1$, alors on a, d'après
une formule de majoration classique pour la
somme des puissances $p$\textsuperscript{ièmes} des inverses des
modules des zéros d'une fonction entière
\[
f(z)=\det(1-zA)\qquad(0<p<1)
\]
\[
s_p(A)\leq p^2\int_0^\infty\frac{\log M_A(r)}{r^{p+1}}\,dr\tag*{\deleted{(10)}}
\]
or $M_A(\unsure{v|r})\leq M_{|A|}(\unsure{v|r})$ (\unsure{\ldots\ --\ --}) d'où
\[
s_p(A)\leq p^2\int_0^\infty\frac{\log M_{|A|}(\unsure{v|r})}{{\unsure{v|r}}^{\unsure{p-1|p+1}}}\,dr
\]
\,; appliquant alors la formule
(9), on obtient donc la majoration
\[
s_p(A)\leq\frac{p\pi}{\sin p\pi}S_p(A)\tag*{(10)}
\]
Le facteur $\frac{p\pi}{\sin p\pi}$ est $>1$, mais tend vers 1
pour $p\to0$. Soit alors $0<p\unsure{\leq|<}+\infty$ quelconque,
et soit $n$ un entier $>0$ tel que $\frac{p}{n}<1$.
On a évidemment $s_p(A)=s_{p/n}(A^n)$, d'où, en
vertu de (10) :
\[
s_p(A)\leq\frac{\frac{p}{n}\pi}{\sin\frac{p}{n}\pi}S_{p/n}(A^n)
\]
Or, on a \deleted{\unsure{$S_{p/n}(A^n)\leq S_p(A)$}}
\[
\|A^n\|_{p/n}\leq\underbrace{\|A\|_p\ldots\|A\|_p}_{n\ \text{facteurs}},
\]
comme il résulte aussitôt de la formule (6), qui sera.
démontrée plus bas (alors il est immédiat
que cette formule se généralise au produit d'un
nb fini quelconque d'opérateurs), d'où
\[
S_{p/n}(A^{\unsure{{-}|n}})=\left(\|A^n\|_{p/n}\right)^{p/n}
\leq\|A\|_p^p=S_p(A).
\]

\scanpage{133}
3). La preuve de la formule (6) est la plus difficile.
Nous commençons par donner une
preuve directe et très simple pour le cas où
\deleted{$\unsure{p,q,r\geq1,\ r=1}$} \inserted{$p\geq1,q\geq1,r=1$} *
(preuve qui a d'ailleurs l'avantage
de \inserted{devoir} s'appliquer sans modification au cas d'une
trace \unsure{$\tau$|$\sigma$|$\operatorname{Tr}$|$D$} au sens de von Neumann - Dixmier, définie
sur une algèbre auto-adjointe $\unsure{\mathcal A|\mathfrak A}$ \deleted{\unsure{ft fermée}}
\deleted{\unsure{unift fermée}} unif\textsuperscript{t} fermée \unsure{fortement}
(1) (\deleted{\unsure{les éléments de}} \textsuperscript{(1)} les ``\deleted{\unsure{éléments}} \inserted{opérateurs} de puissance
$p$\textsuperscript{ième} sommable'' de $\unsure{\mathcal A|\mathfrak A}$ \inserted{-- au sens de la trace $\unsure{\tau|\sigma|\operatorname{Tr}|D}$ --}
se définissant encore
comme l'ens des $A\in\unsure{\mathcal A|\mathfrak A}$ tels que
$\unsure{\tau|\sigma|\operatorname{Tr}|D}(|A|^p)<+\infty$.)\textsuperscript{\unsure{*|a|A}}.
\margin{?? Il n'est d'ailleurs pas difficile \deleted{\unsure{de}} étendre alors le même \deleted{\unsure{tout}} énoncé au cas où $r\geq1$}
(Au N° suivant, nous
obtiendrons d'ailleurs le $=$ résultat comme
conséquence d'une proposition plus générale).

Nous commençons par prouver que, si
$A_1,\ldots A_{2^n}$ sont $2^n$ opérateurs cpt continus,
($n$ entier $\geq1$) alors on a \textsuperscript{(1)}
\[
\|A_1 A_2\mathinner{\cdotp\cdotp\cdotp\cdotp} A_{2^n}\|_1\leq\|A_1\|_{2^n}\cdot\|A_2\|_{2^n}\ldots\|A_{2^n}\|_{2^n}\tag*{(11)}
\]
Nous le prouvons par récurrence sur $n$\,;
c'est vrai si $n=1$, \deleted{\unsure{on a bien}} \inserted{(formule\quad)}, supposons
\deleted{\unsure{$\|A_1A_2\|_1\leq\|A_1\|_2$}}
-la démontrée \inserted{\unsure{la|le} formule analogue} pour $2^{n-1}$ facteurs, et démontrons
alors (11). Écrivant $A_1\ldots A_{2^n}=(A_1A_2)\ldots(A_{2^n-1}A_{2^n})$
le premier membre sera, en vertu de l'hyp.
de récurrence, majoré par $\|A_1A_2\|_{2^{n-1}}\ldots\|A_{2^n-1}A_{2^n}\|_{2^{\unsure{n-1}}}$

\noindent\rule{\linewidth}{0.3pt}

(1) Le lecteur notera que ce résultat suffirait pour
démontrer \deleted{\unsure{directement|exactement}} \inserted{directement}, comme nous le ferons plus bas,
la formule \unsure{(6)|(6')} dans le cas général.

\scanpage{134}
de sorte qu'on est ramené à démontrer \unsure{des|de}
\unsure{formules|formule} \unsure{telles|telle} que
$\|A_1\unsure{A_1|B_1|A_2}\|_{2^{n-1}}\leq\|A_1\deleted{\unsure{\|_{2^2}|A_2}}\|_{2^n}\|\unsure{A_2|B_1}\|_{2^n}$, soit, précisément,
\uline{$\|AB\|_{2^{n-1}}<\|A\|_{2^{\unsure{n|n-1}}}\,\|B\|_{2^n}$.}
Or le premier membre s'écrit
\[
\|AB\|_{2^{n-1}}=\left(\bigl\|\bigl(AB(AB)^*\bigr]^{\frac{2^{n-1}}{2}}\bigr\|_1\right)^{\frac1{2^{n-1}}}
=\left(\|C\|_1\right)^{\frac1{2^{n-1}}}
\]
où
\[
\unsure{C|C_1|C_n}=(ABB^*A^*)^{2^{n-2}}=A\,(BB^*)\,(A^*A)\ldots(BB^*)\,A^*
\]
($2^{n-1}+1$ facteurs au dernier membre) on a.
\[
\|C\|_1=\operatorname{Tr} C\deleted{\unsure{=\operatorname{Tr}}}
\inserted{=\operatorname{Tr}\unsure{\deleted{C}C'|\deleted{H}C'|C'}}\quad
\inserted{\text{où }C'=}\,A^*A\,.\,BB^*\,.\,A^*A\,.\,\ldots\,BB^*
\deleted{\unsure{)=\operatorname{Tr}H}}
\]
\deleted{\unsure{où dans le dernier membre} \unsure{figurent|figurant} $2^{n-1}$}
\unsure{\uline{facteurs}|\deleted{facteurs}}
($2^{n-1}$ facteurs au second membre)

Or $|\operatorname{Tr} C'|\leq\|C'\|_1$, et d'après l'hypothèse de récurrence, on a de nouveau
\margin{\unsure{L|l}}
\[
\begin{aligned}
\|C'\|_1&\leq\|A^*A\|_{2^{n-1}}\|BB^*\|_{2^{n-1}}\ldots\|BB^*\|_{2^{n-1}}\\
&=\bigl(\|A^*A\|_{2^{n-1}}\bigr)^{2^{n-2}}\bigl(\|B^*B\|_{2^{n-1}}\bigr)^{2^{n-2}}
\end{aligned}
\]
Or $\|A^*A\|_{2^{n-1}}=\bigl(\|A\deleted{\unsure{_2|_{2^n}}}\|_{2^n}\bigr)^2$,
$\|B^*B\|_{2^{n-1}}=\bigl(\|B\|_{2^n}\bigr)^2$
d'où $\|C'\|_1\leq\bigl(\|A\|_{2^n}\|B\|_{2^n}\bigr)^{2^{n-1}}$. On obtient donc bien.
\[
\|AB\|_{2^{n-1}}=\bigl(\|C\|_1\bigr)^{\frac{\unsure{2^1|\deleted{2}\inserted{1}}}{2^{n-1}}}
\leq\bigl(\|C'\|_1\bigr)^{\frac1{2^{n-1}}}\leq\|A\|_{2^n}\|B\|_{2^n},
\]
ce qui prouve la formule (11)

\medskip\noindent\hspace*{2em}-\enspace-\enspace-

\scanpage{135}
\deleted{\unsure{2)|4)}} Soient $p\geq\unsure{1|0}$,
$q\leq0$ \quad $\frac1p+\frac1{\unsure{q}}=1$,
$A\in\mathcal L^p$, $B\in\mathcal L^q$, alors $AB\in\mathcal L^1$, et
$\|AB\|_1\leq\|A\|_p\,\|B\|_{\unsure{q|r}}$

\inserted{\uline{Preuve} \unsure{On} \unsure{peut} \unsure{supposer} $A,B\in\mathcal L^1$}
On a $A=UH$, $B=KV$ ($U,V$ partiellement \unsure{isométriques})

Pour \deleted{\unsure{tout|tt}} $n$, soient $\alpha,\beta$ deux entiers tels que $\alpha+\beta=2^n$;
posons $h=H^{\frac1\alpha}$, $k=K^{\frac1\beta}$, \ldots
\[
AB=Uh^\alpha k^\beta V\qquad\text{d'où}\quad\|AB\|_1\leq\|h^\alpha k^\beta\|_1\quad\leq
\]
or (\unsure{Lem.|Lemm.|Lemme} 1)
\[
\begin{aligned}
\|h^\alpha k^\beta\|_1&\leq
\underbrace{\|h\|_{2^n}\ldots\|h\|_{2^n}}_{\alpha}
\underbrace{\|k\|_{\unsure{2^n|p}}\ldots\|\unsure{k|h}\|_{2^n}}_{\beta}\\
&=\bigl(\operatorname{Tr} h^{2^n}\bigr)^{\frac\alpha{2^n}}
\bigl(\operatorname{Tr}(k^{2^n})\bigr)^{\frac\beta{2^n}}\qquad\text{d'où}
\end{aligned}
\]
\[
\|AB\|_1\leq\bigl(\operatorname{Tr}H^{\deleted{\unsure{2|\frac{2^n}{\alpha}}}\lambda}\bigr)^{\frac1\lambda}
\bigl(\operatorname{Tr}K^\mu\bigr)^{\frac1{\unsure{r|\mu|p}}}
\quad\text{où }\lambda=\frac{2^n}{\alpha},\ \mu\unsure{=|:}\frac{2^n}{\beta}
\]
\deleted{\unsure{Faisons} alors}

Cette inégalité est vraie pour \inserted{un couple} $(\lambda,\mu>1$ pourvu qu'on puisse les mettre sous la \deleted{\unsure{forme}} forme
$\frac{2^n}{\alpha}\unsure{,|\cdot}\frac{2^n}{\beta}$ avec $\alpha\unsure{+|-}\beta=2^n$,
i.e pourvu $\frac1\lambda+\frac1\mu=1$, \deleted{\unsure{que|et|$\lambda p$}}
et $\frac1\lambda$ et $\frac1\mu$ sont des multiples de $\frac1{2^n}$;
par un passage à la limite immédiat, on conclut
\[
\|AB\|_1\leq\bigl(\operatorname{Tr}H^p\bigr)^{\frac1p}
\bigl(\operatorname{Tr}H^{q'}\bigr)^{\frac1{p'}}
\]
\inserted{\unsure{On|la} \unsure{démontre|démonstr.}, de =, si $\frac1{p_1}+\frac1{p_{\unsure{2}}}\ldots\unsure{+|-}\frac1{p_n}=1\unsure{)|,}$
$\|A_1\ldots A_n\|_1\leq\|A_1\|_{p_1}\ldots\|A_n\|_{p_{\unsure{n}}}$}

3) On en conclut \unsure{qu'on|que l'on} a ($p\geq1,\ \frac1p+\frac1{p'}\unsure{=|\geq|\leq}1$)
\[
\|A\|_p=\operatorname*{Sup}_{\|B\|_q\leq1}|\operatorname{Tr} AB|\qquad\text{(d'où \unsure{facilement}}
\]
\[
\|A_1+A_2\|_p\leq\|A_1\|_p+\|A_2\|_p\quad)\quad\text{et enfin}
\]
\deleted{$AB$} si $\frac1p+\frac1q\geq\frac1r$ \quad ($p\unsure{\geq|>}1,q\geq1,\unsure{r}\geq1$)

\scanpage{136}
Le résultat précédent nous permet de démontrer le résultat auxiliaire suivant, qui nous permettra de démontrer la formule (6) en \deleted{\unsure{toute}} généralité :

\uline{Lemme} Soient $A$ et $B$ deux opérateurs $\in\mathcal L^{2m}$ \unsure{($\mathcal H_2$)|($\mathcal H$)|($\mathfrak H_2$)|($\mathfrak H$)|($H_2$)|($H$)},
\margin{posons\;\rule[-2.1em]{0.4pt}{3.3em}}
$H=A^*A,\ K=B^*B,\ L=(AB)^*(AB)$.
Soient \deleted{\unsure{$L$ soit|et soit}} $m$ et $n$ des entiers $\geq1$.
Alors \deleted{\unsure{$L^m$ est, $H$|$L^n,H^n$|$L^m\in\mathcal L^1,H$}} \inserted{$L^m,H^m$} et $K^m$ sont opérateurs à trace, et on a
\[
\alpha_n\bigl(|\deleted{\unsure{A}}L|^m\bigr)
\leq\alpha_n\bigl(\deleted{\unsure{H}}H^m\bigr)\,\alpha_n\bigl(\unsure{\deleted{\unsure{K}}K^m|\lvert K^m}\bigr)
\tag*{(12)}
\]
Le fait que $L^m,H^m,K^m$ soient opérateurs à trace résulte \deleted{\unsure{enfin d'un des précédents}} \inserted{aussitôt de ce qui précède}. (12) s'écrit alors, \deleted{\unsure{cela|et|entiers}}
\deleted{\unsure{par contre s'écrit en effet en posant}}
\inserted{en posant} $\mathscr A=\underline{\wedge} A,\ \mathscr B=\underline{\wedge} B$
\inserted{$\mathscr H=\mathscr A^*\mathscr A,\ \mathscr K=\mathscr B^*\mathscr B,\ \mathscr L=(\mathscr A\mathscr B)^*(\mathscr A\mathscr B)$}
\[
\operatorname{Tr}\,|\deleted{\unsure{\mathscr A\mathscr B}}\mathscr L|^{\unsure{m|2m}}\leq
\operatorname{Tr}\,\deleted{\unsure{\mathscr H}}\mathscr H^m\ \operatorname{Tr}\,\deleted{\unsure{\mathscr B}}\unsure{\lvert\mathscr K\rvert^m|\mathscr K^m},
\]
\deleted{\illegible{2}}
\[
\deleted{\unsure{\bigl(\|\mathscr A\mathscr B\|_m\bigr)^m\leq
\bigl(\|\mathscr A\|_m\bigr)^m\bigl(\|\mathscr B\|_m\bigr)^m|\bigl(\|L\|_m\bigr)^m\leq\bigl(\|H\|_m\bigr)^m\bigl(\|K\|_m\bigr)^m}}\quad
\deleted{\text{\unsure{ce qui}}}
\]
\deleted{$\|\mathscr A\mathscr B\|_m\leq\|\mathscr A\|_m\|\mathscr B\|_m$. Or, \unsure{on a vu qu'on a même}
$\|AB\|_m\unsure{<|\leq}\|A\|_{\frac m2}\|B\|_{\frac m2}$}

ce qui nous ramène au cas où $n=1$. On doit alors démontrer
\[
\unsure{\operatorname{Tr}\||\operatorname{Tr}\lvert|\operatorname{Tr}[}\,(ABB^*A^*)\mathinner{\cdotp\cdotp\cdotp-}(ABB^*A^*)\|_2
\leq\bigl(\|\deleted{\unsure{H}}A^*A\|_m\bigr)^m\bigl(\|B^*B\|_m\bigr)^m
\]
Or, le premier membre \deleted{\unsure{est la trace-norme}} \inserted{est aussi égal à la}
\deleted{\unsure{d'un produit de $4m$ facteurs, est en effet}} trace du produit de $2m$ facteurs
\deleted{\unsure{majoré par}}
\[
(A^*A)(BB^*)\ldots(BB^*),
\]
donc au plus égal \unsure{à|:} la trace-norme de ce produit, donc au plus égal \unsure{à|:}
$\bigl(\|A^*A\|_{2m}\bigr)^m\bigl(\|BB^*\|_{2m}\bigr)^m$
d'après la formule $\|A_1\ldots A_{2m}\|_1\leq\|A_1\|_{2m}\ldots\|A_{2m}\|_{\unsure{2m}}$

Or, on a $\|A^*A\|_{2m}\leq\|A^*A\|_m$, et
$\|BB^*\|_{2m}\leq\|B^*B\|_{2m}\leq\|B^*B\|_m$, d'où résulte le lemme.

\scanpage{137}
Soient maintenant \inserted{$0<p,q,r\leq\infty$} $A\in L^p$ $A\in L^q$, $\frac1p+\frac1q\geq\frac1r$ démontrons la formule (6). Nous pouvons supposer \deleted{\unsure{$p,$|$r$}} que $p$ et $q$ \deleted{\illegible{2}}\inserted{sont finis} \deleted{\unsure{puisque deux infinis} (\unsure{on en aura alors,}} \inserted{\unsure{puisque si l'un était infini,} les formules (2)} donneraient le résultat; \deleted{$\|AB\|_{\unsure{r|\infty}}$ ($\|AB\|_\infty=\|A\|_\infty\|B\|_\infty$) \unsure{est un|et on}} \inserted{\unsure{de|De} plus,} peut supposer évidemment \deleted{\unsure{ici|$\frac1r$|$\frac1q$}} que l'on a exactement $\frac1p+\frac1q=\frac1r$ (puisque $\|X\|_r$ est fonction décroissante \unsure{de|en} $r$)

Soit $m$ \unsure{un|\deleted{un}} entier tel que $p\leq 2m$, $q\leq 2m$. \deleted{\unsure{Alors|Pour} $A$} Posons \unsure{comme} dans le lemme 2

\deleted{On a}

\[
\deleted{S_r(AB)\unsure{\leq|=} S_{\frac r2}(AB(AB)^*)=S_{\frac r{2m}}\unsure{A|L}}
\]

$L=(AB)^*(AB)^{\deleted{\unsure{*}}}$, $H=A^*A$, $K=B^*B$, on aura.

\[
S_r(AB)=S_{\frac r2}(L)=S_{\frac r{2m}}(L^m)\deleted{\unsure{\leq p+2m|=r/2m}}
\]

d'où, en vertu de la formule (\unsure{9|1}), \unsure{en posant|on peut} \unsure{pour} la commodité $\lambda(x)=\frac{\pi}{\sin\pi x}$

\[
S_r(AB)=S_{\frac r{2m}}(L^m)=\frac{r/2m}{\lambda(r/2m)}
\int_0^\infty\frac{\log M_{L^m}(\rho)}{\rho^{1+\frac r{2m}}}\,d\rho \tag*{(11)}
\]

\deleted{\unsure{Or, on a} $\unsure{M_{L^m}|M}$} Or, $M_{L^m}(\rho)=\sum\alpha_n(L^m)\rho^n$, et \deleted{$M_{L^m}(\rho)=\sum_n\alpha_{\unsure{n|k}}(L^{\unsure{m|k}})\rho^{\unsure{k|n}}$} \inserted{$\alpha_n(L^m)\leq\alpha_n(H^m)\alpha_n(K^m)$} en vertu du lemme 2 on a \deleted{$\alpha_n(L^m)\leq\alpha_n(H^{\unsure{2|m}})\unsure{\alpha_n K|\alpha_n H}$}

\margin{\uline{\unsure{Ici lemme|Le lemme|Ici le cours}} sur \uline{les fon}ctions \uline{entiè}res}

d'où

\[
M_{L^m}(\rho)\leq\sum\alpha_n(H^m)\alpha_n(K^m)\rho^n
\]

Posons $\left[p'=\frac1p\;\;q'=\frac1q\;\;r'=\frac1r\right]$, on : \deleted{$\unsure{p'+q'=r'}$}

\[
\inserted{p'r+q'r=1}
\]

d'où
\margin{\tikz{\draw[line width=0.3pt] (0,0) -- (0.6,0.1) .. controls (0.55,0.35) and (0.46,0.65) .. (0.46,0.9);\draw[line width=0.3pt] (0.5,0.95) -- (1.35,1);}}

\[
\begin{aligned}
M_{L^m}\unsure{(}\rho)&\leq\sum\alpha_n(H^m)\alpha_n(K^m)\rho^{np'r+nq'r}\\
&=\sum_n\bigl(\alpha_n(H^m)\rho^{np'r}\bigr)
\bigl(\alpha_n(K^m)\rho^{nq'r}\bigr)
\end{aligned}
\]

\scanpage{138}
\[
\begin{aligned}
&\leq\left(\sum_n\alpha_n(H^m)\rho^{np'r}\right)
\left(\sum_n\alpha_n(K^m)\rho^{nq'r}\right)\\
&\deleted{\unsure{=|\leq}}\qquad=M_{H^m}(\rho^{p'r})M_H(\rho^{q'r})
\end{aligned}
\]

d'où, compte tenu de (12)

\[
\frac{\lambda(r/2m)}{r/2m}S_r(AB)\leq
\int_0^\infty\frac{\log M_{H^m}(\rho^{p'r})}{\rho^{1+r/2m}}\,d\rho+
\int_0^\infty\frac{\log M_{K^m}(\rho^{q'r})}{\rho^{1+r/2m}}\,d\unsure{\rho}
\tag*{(13)}
\]

Calculons \deleted{\unsure{par exemple}} la première intégrale du second membre, par le changement de variable $\sigma=\rho^{p'r}$ d'où $\rho=\sigma^{pr'}$ $\frac{d\rho}{\rho}=pr'\frac{d\sigma}{\sigma}$ il vient

\[
\int_0^\infty\frac{\log M_{H^m}(\rho^{p'r})}{\rho^{r/2m}}\frac{d\rho}{\rho}
=pr'\int_0^\infty\frac{\log M_{H^m}(\sigma)}{\sigma^{p/2m}}\frac{d\sigma}{\sigma}
\]

Le deuxième membre est égal (\unsure{\deleted{formule}|formule} (9)) à

\[
pr'\frac{\lambda(p/2m)}{p/2m}S_{p/2m}(H^m)=\frac{\lambda(\unsure{p|p/2m})}{r/2m}S_p(A).
\]

\unsure{De} $=$ la deuxième intégrale dans (13) est égale à $\frac{\lambda(q/2m)}{r/2m}S_q(B)$, et \deleted{\unsure{on obtient|(13) est}} \inserted{(13) s'écrit \unsure{ainsi}}

\[
\lambda(r/2m)S_r(AB)\leq\lambda(p/2m)S_p(A)+\lambda(q/2m)S_q(B)
\]

ou aussi

\[
S_r(AB)\leq\frac{\sin\pi\frac{r}{2m}}{\sin\pi\frac{p}{2m}}S_p(A)+
\frac{\sin\pi\frac{r}{2m}}{\sin\pi\frac{q}{2m}}S_q(B)
\]

\inserted{Cela prouve déjà que $S_{\unsure{r|p}}(AB)\leq+\infty$, et donne même une majoration effective de $S_r(AB)$ en fonction de $S_p(A)$ et $S_q(B)$.}

Faisant tendre $m$ vers $+\infty$, on obtient la formule plus simple

\[
\boxed{S_r(AB)\leq\frac{r}{\deleted{2m}p}S_p(A)+\frac rq S_q(B)}\tag*{(14)}
\]

C'est une inégalité remarquable, mais n'est pas celle que nous cherchons. Mais \deleted{\unsure{cette relation} \illegible{1}}\textsuperscript{\deleted{\illegible{1}}}

\scanpage{139}
\deleted{\unsure{du fait d'homogénéité|à part d'homogénéité} \unsure{il vient,|évidente,} qui permet}

Remplaçons dans cette formule $A$ par $\lambda A$, $B$ par $\mu B$ ($\lambda$ et $\mu$ nombres $>0$) on obtient.

\[
\deleted{\unsure{\lambda^p\mu^q=\frac{\lambda^p\mu^q}{r}S_r(AB)}}
\inserted{\frac{(\lambda\mu)^r}{r}S_r(AB)}
\leq\frac{\lambda^p}{p}S_p(A)+\frac{\mu^q}{q}S_q(B)
\]

\deleted{Exprimons que cette inégalité doit être \unsure{vérifiée} quels que soient $x,y\unsure{\geq|>}0$.}

\margin{\unsure{Ici lemme que} la condition du \unsure{n°} implique $c^{\frac1r}\leq a^{\frac1p}b^{\frac1q}$ i.e \uuline{$c\leq a^\alpha b^\beta$}}

Pour simplifier, posons $a=S_p(A)$, $b=S_q(B)$, $c=S_{\unsure{p|r}}(C)$, on \unsure{aura}

\[
c\leq\mathop{\mathrm{Inf}}\limits_{\substack{\lambda\geq0\\\mu\geq0}}
\frac{a\frac{\lambda^p}{p}+b\frac{\mu^q}{q}}
{\substack{\deleted{\unsure{(\lambda\mu)^r}}\\\deleted{\unsure{\frac{\lambda^p\mu^q}{r}}}}\frac{(\lambda\mu)^r}{r}}
\]

\deleted{Posons $\unsure{x^p=\lambda,y^q}$}\deleted{\inserted{$x=\lambda^p,\ y=\mu^q$,}}

Posons $x=\lambda^p$, $y=\mu^q$, $\alpha=\frac rp$, $\beta=\frac rq$ \unsure{.|,} d'où \unsure{,} \uline{$\frac rp+\frac rq=1$},

On a
\[
c\leq\mathop{\mathrm{Inf}}\limits_{\substack{x>0\\y>0}}
\frac{\alpha a\deleted{\unsure{\frac rp}}x+\deleted{\unsure{\frac rq}}\beta by}{x^\alpha y^\beta}
=\mathop{\mathrm{Inf}}\limits_{\substack{x\geq0\\y\geq0\\x^\alpha y^\beta=1}}\alpha a x+\beta by
\]

Or on vérifie aussitôt que \unsure{la|le} \deleted{\unsure{minimum}}\inserted{borne \unsure{inférieure}} du second membre est atteinte (car l'en\unsure{s.} \unsure{des} points $(x\unsure{,|;}y)$ tels que $x\geq0,y\geq0$ $ax+by\leq M$ -- valeur donnée -- est compact) elle est donc atteinte en un point $(x,y)$ tel que les équations \deleted{$\unsure{a\,dx+b\,dy=0}$}\inserted{$d(\alpha ax\unsure{-|+}\beta by)=0$}, $d(x^\alpha y^\beta)=0$ \deleted{soient compatibles} \inserted{aient une solution non nulle en} $dx,dy$ ;

\unsure{or} $d(x^\alpha y^\beta)=\alpha x^{\alpha-1}y^\beta\,dx+\beta y^{\beta-1}x^\alpha\,dy$,\margin{\unsure{2|?}}
\deleted{i.e\unsure{.} tel que $\unsure{\det(\alpha a\quad\beta b\quad\alpha a}$}
la condition d'extremum est donc $\frac{\alpha a}{\alpha x^{\alpha-1}y}=\frac{\beta b}{\beta x^\alpha y^{\beta-1}}$

soit $ax=by$ \deleted{\unsure{$=$}}, \deleted{\unsure{et comme on a}} \inserted{d'où $ax=by=\frac{\alpha ax+\beta by}{\alpha+\beta}$} ; \deleted{$\unsure{\alpha ax+by=\frac{\alpha ax}{\unsure{\alpha|2}}\unsure{=|+}\frac{\alpha ax+\beta by}{\unsure{\alpha|2}}}$} Comme $\alpha+\beta=1$, le minimum cherché sera alors \deleted{\unsure{$ax=by$|$ax+by$|$\alpha x+b$}} égal à $ax=by$ ; compte tenu de la \unsure{condi}

\scanpage{142}
\uline{Proposition} Pour que la suite $(H_i)$ d'él. hermitiens $\in L^p(\mathscr H)$ converge vers $H_0\in L^p(\mathscr H)$, il faut et il suffit que

\[
\text{1°)}\quad\|H_i\|_p\to\|H_0\|_p\qquad
\text{(ou seulement }\limsup_{i\to\infty}\|H_i\|_p\leq\unsure{\|H_0\|_p|\lvert H_0\|_p}
\]

2°) Pour \deleted{tout} $\lambda$ du spectre de $H_0$, (soit $\alpha>0$ tel qu'il \deleted{$[\lambda-\alpha,\lambda+\alpha]$ ne \unsure{contien|contienne}}\inserted{n'existe pas de} valeur spectrale $\mu\neq\lambda$ de $H_0$ telle que $|\lambda\deleted{\unsure{+}}-\mu|\leq\zeta$, \deleted{\unsure{élément spectral de $H$|l'ensemble spectral de $H$} \unsure{faut que}} \deleted{et telle} et tel que $\lambda-\alpha$ et $\lambda+\alpha$ n'appartiennent au spectre \deleted{\unsure{de} $H_{\unsure{0|i}}$}\inserted{$h_i$}\inserted{d'aucun $H_i$}) \deleted{\unsure{alors}} la composante spectrale de $H_i$ \deleted{\unsure{relative à}} dans $[\lambda-\alpha,\lambda+\alpha]$ est de dimension \deleted{\unsure{uniforme égale à}} qui finit par être égale à celle de $H_0$, et $h_i\to h_0$ (inutile de spécifier la topologie, en vertu d'une proposition préliminaire sur l'identité de \deleted{toutes} ces topologies)

\margin{\unsure{évidem\textsuperscript{t}|évident}\\ faux ds\\ le cas non\\ hermitien}

\uline{Corollaire 1} Les deux conditions suivantes sont néc. et suffis. \unsure{stes} pour la convergence de la suite $H_i$ vers $H_0$

\[
\text{1)}\quad\|H_i\|_p\to\|H_0\|_p\qquad\text{(resp$^{\unsure{\mathrm t|\mathrm d}}$ .\ .\ --)}
\]

\[
\text{2)}\quad H_i\to H_0\quad\text{dans }L^\infty(\mathscr H)
\]

\uline{Corollaire 2} $A\to |A|^r$ est application continue de $L^p$ dans $L^{\frac pr}$.

En effet, elle est composée de $A\to A^*A$, continue de $L^p$ dans $L^{\frac p2}$, avec

\scanpage{143}
$H\to H^{\frac r2}$ de $L^{+\;\frac p2}$ dans $L^{\frac pr}$ ; or, les conditions de la prop. 1 sont aussitôt vérifiées \unsure{(|\deleted{(}}\deleted{\unsure{en effet}} (condition 1 triviale ; condition 2 se voit \deleted{\unsure{ainsi}} en se ramenant à \deleted{\unsure{montrer que si}}\inserted{la dimension} finie, \unsure{où|ou|on} c'est trivial ($\sum_{i=\unsure{0|1}}^{\unsure{n|-}}\lambda_i a_i\otimes a_i$, $\lambda_i\geq0$ et $a_i$ variables, $\to\lambda_i^{(0)}$ et $a_i^{(0)}$ ; \unsure{on|en} \unsure{forme} $\sum\lambda_i^{r/2}a_i\otimes a_i$ . --- --)

\uline{Remarque} 1 Si $0<\unsure{p|r}\leq2$, pour qu'une suite d'hermitiens $H_i$ dans $L^p$ converge \inserted{\unsure{vers} $H_0$}, il faut et il suffit qu'elle converge faiblement, et que $\|K_i\|\to\|H_0\|$ (ceci est connu au cas hilbertien)

\margin{\deleted{\uline{\unsure{prendre comme}} \uline{\unsure{indépendant de $p$|pertinent ds $L^p$}} \unsure{la norme uniforme de Murray|On connaît le contre-exemple de Murray}}}\margin{!}\margin{\unsure{vu}}

Est-ce vrai si $p<+\infty$ ? Uniforme convexité de $L^p$ ?

\uline{Remarque 2} Est-ce que $S_p\unsure{\lvert|\lVert|(\lvert} H^r-K^r\unsure{\rVert|\rvert}^{\frac1r}$ est fonction croissante de $r$ ($H$ hermitien positif) ? -- Il suffirait de \deleted{\unsure{le|ce}} savoir que $r<1$ \unsure{implique|impliq.}

\[
\bigl\|\lvert H^r-K^r\rvert^{\frac1r}\bigr\|\leq\|H-\unsure{K|k}\|
\]

--- Alors on aurait une précision notable dans le corollaire 1.

\scanpage{144}
\margin{\begin{tikzcd}[ampersand replacement=\&,column sep=tiny]
E\arrow[r,"u"]\&F\arrow[r,"v"]\&E
\end{tikzcd}}
\margin{\uline{majorations individuelles}}

\noindent\(\left\{\begin{minipage}{0.94\linewidth}
\[
\left\{\begin{aligned}
\rho_{m+n-1}(vu)&\leq\rho_m(\unsure{v|u})\rho_n(u)\\
\rho_{m+n-1}(u+v)&\leq\rho_m(u)+\rho_n(v)
\end{aligned}\right.
\]

En effet,

\[
\rho_{m+n-1}(vu)=\mathop{\mathrm{Inf}}_{E_{m+n-\unsure{2|1}}}
\mathop{\mathrm{Sup}}_{\ni x,\,\|x\|\leq1}\|vu x\|
\]

Or il existe $E_{n-1}\subset E$ tel que $\mathop{\mathrm{Sup}}_{\substack{\|x\|\unsure{\leq|<}1\\x\in E_{n-1}}}\deleted{\unsure{\rho|E|\Sigma}}\|ux\|=\rho_n(u)$ \deleted{\unsure{et|est}}, il existe $F_{m-1}\subset F$ tel que $\mathop{\mathrm{Sup}}_{\substack{\|y\|\leq1\\y\in F_{m-1}}}\|vy\|=\rho_m(v)$ soit alors \deleted{$\unsure{E_{m-2}|E_{m+n-2}}$}$\inserted{H}=u^{-1}(F_{m-1})$, on a

\[
\mathop{\mathrm{Sup}}_{\substack{x\in H\unsure{\cap|\wedge} E_{n-1}\\\|x\|\leq1}}\|vux\|\leq\rho_n(u)\rho_m(v)
\]

or $p-\deleted{\unsure{p}}1=\operatorname{codim}H\unsure{\cap|\wedge} E_{n-1}\geq\operatorname{cod}H+\operatorname{cod}E_{n-1}\geq m-1+n-1=(m+n-1)-1$

\unsure{donc|d'où}

\[
\rho_p(vu)=\mathop{\mathrm{Inf}}\mathop{\mathrm{Sup}}_{\substack{x\unsure{\perp|\in} E_{p-1}\\\|x\|\leq1}}\|vux\|\leq\rho_n(u)\rho_m(v)
\]

à fortiori

\[
\boxed{\rho_{m+n-1}\unsure{(}vu)\leq\rho_n(u)\mathbin{.}\rho_m(v)}
\]

De $=$, si $u,v\in L_0(E,F)$, on a

\[
p-\deleted{\unsure{p}}1=\operatorname{codim}(E_{n-1}, E'_{m-1})\geq m-1+n-1=(m+n-1),
\]

\[
\deleted{\unsure{\rho_p(u+v)=\mathrm I}}
\mathop{\unsure{\mathrm I|\mathrm{Inf}}}\mathop{\mathrm{Sup}}_{\substack{x\unsure{\perp|\in} E_{n-1}\unsure{\cap|\wedge} E'_{m-1}\\\|x\|\leq1}}
\|(u+v)x\|\leq\rho_n(u)+\rho_m(v)
\]

d'où à fortiori

\[
\rho_{m+n-1}(u+v)\leq\rho_p(u+v)\leq\rho_n(u)+\rho_m(v\unsure{)}
\]
\end{minipage}\right.\)

\par\noindent\rule{\linewidth}{0.4pt}\par

\unsure{\deleted{C'est}|\uline{C'est}} On peut avoir $\limsup\frac{|\lambda_i(u)|}{\rho_{\unsure{i|n}}(u)}=+\infty$

Majorations individuelles pour $|\lambda_i(u)|\leq\rho_{\unsure{i/2|i/k}}(u)$ ou analogue ?

\par\noindent\rule{\linewidth}{0.4pt}\par

\scanpage{145}
Si $E$ est un Banach, et $u_\alpha\to u$ \unsure{dans|ds} $L_0(E)$, alors $\sigma(u_\alpha)\to\sigma(u)$ dans le sens suivant\unsure{.|:}

De plus, convergence des projecteurs spectraux \unsure{--\,--\,--\,--\,--\, .} Dans le cas des op. hermitiens dans l'espa\unsure{ce} de Hilbert, on a $=$ \unsure{là|la} des conditions nécessaires et suffisantes de convergence.

Convergence de \unsure{l'hermitien|d'hermitiens|l'hermiticité} \inserted{$u_\alpha\to u$} dans $L^p(E)$ : il est nécessaire et suffisant que \textcircled{\scriptsize 1°} $\|u_\alpha\|_p\to\|u\|_p$, \textcircled{\scriptsize 2°} $u_\alpha\to u$ spectralement, i.e, dans $L(E)$

(il suffit \unsure{$=$|même} que l'on ait la convergence spectrale individuelle en chaque $\lambda\in\sigma(u)$, $\lambda\neq0$)

\uline{Corollaire 1} $\unsure{U|u}\to|A|$, de $L^p(E,F)$ dans $L^p(E)$ est continue. Plus généralement, $u\to|u|^r$ de $L^p(E,F)$ dans $L^{p/r}(E)$ est continue.

Car composée de $u\to u^*u$ de $L^p(E,F)$ dans $(L^{p/2}(E))^+$, avec $v\to v^{\frac r2}$ de $(L^{p/2}(E))^+$ dans $L^{\frac pr}(E)$ ; or \deleted{\unsure{la|cette}} \unsure{dernière|dernier} est continue, car elle est continue de $L_0(E)^+$ dans $L_0(E)$ (ce dernier \unsure{point|fait} se vérifie \uline{spectralement} ; c'est sans doute vrai aussi de $L(E)^+$ dans $L(E)$ ? Inégalités précises ?)

\uline{Corollaire 2} Si $0<p\leq\unsure{2|1}$, pour qu'une suite \inserted{d'hermitiens} $u_\alpha\to u$ converge dans $L^p(E\deleted{\unsure{,F|,E}})$, il faut et il suffit que $\|u_\alpha\|_p\to\|u\|_p$, et qu'on ait la convergence faible\unsure{.}

\scanpage{146}
Question : \unsure{Dans quel cas (uniforme convexité} \unsure{les|ces|des} \unsure{espaces|esp.} $L^p(E,F)$ \unsure{?} pour $0<p<\infty$ ? Pas d'importance.

\unsure{2|1}) Pour la continuité de $h\to h^r$, \unsure{on} est ramené au cas $0\leq r<1$ (car c'\unsure{est} \unsure{trivial} pour $r$ entier) Il semble possible que

\[
\bigl\|\lvert {\unsure{h|H}}^r-k^r\rvert\bigr\|^{1/r}\leq\|h-k\|\qquad\text{pour }0\leq r\leq1,
\]

Il revient au $=$ de \unsure{prouver|présumer} que

\[
S_{p/r}(h^r-k^r)=S_p\bigl(\lvert {\unsure{h|H}}^r-k^r\rvert^{1/r}\bigr)
\]

est fonction croissante de $r$, ($h$ et $k$ hermitiens positifs, $p$ quelconque)\unsure{.|,}

On peut faire aussi $p=1$.

\scanpage{147}
\uline{Th} Si $A_i\to A$ \unsure{ds|dans} $L^p(\mathscr H)$, alors $\widehat A_{\unsure{i}}\to\widehat A$ \unsure{ds|dans} $\Sigma_p$
\unsure{Ceci vient|Cela résulte} facilement \unsure{du|un}

\uline{lemme} Pour \unsure{tt|tout} $\varepsilon>0$, on a $\rho>0$ et $i_0$ tel que $i>i_0$, implique

\[
\sum_{r_\nu(A_i)\leq\rho}\bigl(r_\nu\unsure{(A_i)|[A_i]}\bigr)^p\leq\varepsilon
\]

Pour \unsure{s'en|nous} convaincre, soit $\unsure{w|\omega}(t)$ la fonction convexe qui est identique à $e^{pt}$ pour $t\leq\log\rho$, et linéaire pour $t\geq\log\rho$, i.e $\unsure{-}\unsure{w|\omega}=p\rho^p(t-\log\rho)+\rho^p=\rho^p\log e\left(\frac{e^t}{\rho}\right)^{\unsure{p}}$ on a (posant $r_\nu=\unsure{r|v|\nu}_\nu(A_i)$, \deleted{\unsure{$s_\nu$ analogue}}\inserted{$s_\nu=\unsure{r|v|\nu}_\nu(|A_i|)$,})

\margin{\unsure{On n'a pas besoin} de la continuité de $A\to|A|$, si on \unsure{remplace} $|A|$ par $A^*A$\\$\unsure{>}$}

\[
\sum \unsure{w|\omega}(\log\unsure{r|v|\nu}_\nu)\leq\sum \unsure{w|\omega}(\log(s_\nu))
\]

soit

\[
\sum_{r_\nu\leq\rho}r_\nu^p+\rho^p\sum_{r_\nu\unsure{\geq|>}\rho}\log e\left(\frac{r_\nu}{\rho}\right)^p
\leq\sum_{s_\nu\leq\rho}s_\nu^p+\rho^p\sum_{s_\nu\unsure{\geq|>}\rho}\log e\left(\frac{s_\nu}{\rho}\right)^p
\]

On peut donc \uline{supposer les $A_i\geq0$}, \inserted{\unsure{(}car $|A_i|\to|A|$)} et prouver que alors, pour \unsure{tt|$\forall$|tout} $\varepsilon>0$, on a $\rho>0$ et $i_0$ tel que $i>i_0$ implique

\[
\sum_{\unsure{r_\nu(A_i)|r_\nu\lvert A_i\rvert}\leq\rho}\unsure{r_\nu(A_i)^p|r_\nu\lvert A_i\rvert^p}
+\rho^p\sum_{r_\nu(A_i)\unsure{\geq|>}\rho}\log e\left(\frac{r_\nu(A_i)}{\rho}\right)^p\unsure{<|\leq}\varepsilon
\]

On prouve que d'abord, par extraction d'une \deleted{\unsure{suite|Soit} $r_\nu\unsure{=}\,\unsure{r|v|\nu}_\nu(A_{\unsure{i'|i}})$, il suffit de montrer séparément}

\margin{\textcircled{\scriptsize \unsure{si Si|$s_iS_i$}}}

\[
\deleted{\text{a)}\qquad\rho^p\sum_{r_\nu\unsure{\geq|>}\rho}\log e\left(\frac{r_\nu}{\rho}\right)^p\to0\qquad\text{si }\rho\to0}
\]

\deleted{b) \unsure{l'énoncé du lemme|l'inégalité du lemme} pour les hermitiens positifs} $\unsure{\deleted{A_i}|A_i}$.

\deleted{\unsure{Or} \illegible{7}}

\uline{Preuve de a)} Or, on a en effet

\margin{\deleted{\illegible{1}}}

\[
\begin{aligned}
\rho^{\unsure{\nu|p}}\sum_{r_\nu>\rho}\log e\left(\frac{r_\nu}{\rho}\right)^p
&=\rho^{\unsure{\nu|p}}\int_{r>\rho}^{+\infty}-\log e\left(\frac{\unsure{r_\nu|r}}{\rho}\right)^p\,dn(r)\\
&=\rho^p\left[-n(r)\log-\right]_{\rho+0}^{+\infty}
+p\rho^p\int_\rho^{+\infty}\frac{n(r)}{r}\,dr\\
&=\rho^p n(\rho)+p\rho^p\int_\rho^{+\infty}\frac{n(r)}{r}\,dr
\end{aligned}
\]

et il est bien connu, puisque $\sum r_\nu^p<\infty$, que $\rho^p n(\rho)\to0$

\scanpage{148}
et que $\int_0^{+\infty}\frac{n(\unsure{r|v})}{\unsure{r|v}^{1-p}}\,d\unsure{r|v}<+\infty$, à fortiori $\int_{\unsure{1|0}}^{+\infty}\frac{n(\unsure{r|v})}{r}\,d\unsure{r|v}<+\infty$,

donc $p\rho^\alpha\int_\rho^{+\infty}\frac{n(\unsure{r|v})}{r}\,d\unsure{r|v}<p\rho^\alpha\int_0^{+\infty}\frac{n(\unsure{r|v})}{r}\,dr\to0$ si $\rho\to0$.

\par\noindent\unsure{\rule{0.95\linewidth}{0.2pt}}\par

\deleted{pour \unsure{tout|tt} $\varepsilon$, on peut \unsure{supposer} \unsure{aussi} que}

\[
\deleted{\sum\|A_{\deleted{\unsure{i}}\unsure{i+1}}-A_i\|_p<+\infty,\quad
\sum S_p(A_{\unsure{i-1|i+1}}-A)<+\infty}
\]

\deleted{alors $A=\sum A_{i+1}-A_i$, \unsure{d'où|donc} $A\leq\sum|A_{i+1}-A_i|$}

\uline{Preuve de b} On a

\[
\begin{aligned}
\sum_{r_\nu(A_{\unsure{i|n}})\leq\rho}(\unsure{r|v}_\nu(A_{\unsure{i|n}}))^{\unsure{r|v|\nu|p}}&=S_p(A_{\unsure{i|n}})-\sum_{r_\nu(A_{\unsure{i|n}})>\rho}(r_\nu(A_{\unsure{i|n}}))^p\\
&\longrightarrow\deleted{\unsure{\sum|\leq}} S_p(A)-\sum_{r_\nu\unsure{>|\geq}\rho}r_\nu^p=\sum_{r_\nu\unsure{\leq|<}\rho}r_\nu^p
\end{aligned}
\]

il suffit de prendre $\rho$ tel que $\sum_{r_\nu\leq\rho}\unsure{r|v}_\nu^{\unsure{p|r}}<\varepsilon$, puis $i$ assez grand.

\scanpage{149}
$\mathscr H$ Hilbert. Si $u$ est à trace, $\frac1{1+u}=1-u\frac{R(u)}{\det(1+u)}$

\margin{$\begin{gathered}R(u)=\sum_0^\infty P_{\unsure{n|k}}(u)\\=\frac{\det(1+u)}{1+u}\end{gathered}$}

\uline{Théorème} \uline{1) Soit $u$ opérateur \deleted{\unsure{herm.}} à trace normal. Mais}

\margin{\unsure{$*$}}
\[
\unsure{\deleted{?}|\deleted{=}}\quad
\boxed{|P_n(u)|\leq P_n(|u|)\ ;\ |R(u)|\leq R(|u|)}\tag*{\unsure{$(*),\,R)$|$(*)$|$(*),\,({**})$}}
\]

En particulier, si $u$ est $\geq0$, il en est de $=$ de $P_n(u)$ et de $R(u)$

2) Si $\unsure{u|U}\geq0$, alors la suite des termes $\alpha_{n-k}(u)u^k$ dans l'expression de $P_n(u)$, $P_n(u)=\sum_{k=0}^n\alpha_{n-k}(u)u^k$, est une suite \uline{décroissante} d'hermitiens positifs \deleted{\unsure{quand $k$}}.

Donc on a en particulier, pour $u\geq0$

\[
\deleted{\unsure{\alpha_n(u)}}\quad
\boxed{0\leq P_n(u)\leq\alpha_n(u)\mathbin{.}1\qquad\text{d'où}\quad
\deleted{\unsure{R(u)\leq\alpha_n(u)\,1}}\inserted{0\leq R(u)\leq\det(1+u)1}}\tag*{(**)}
\]

\deleted{\unsure{3)}} \uline{Corollaire} Si $u$ est opérateur \deleted{\unsure{de rang}} à trace normal, on a

\[
\left\{\begin{aligned}
\|R(u)\|_\infty&\leq\|R(|u|)\|_\infty\leq\det(1+|u|)\leq e^{\|u\|_1}\\
\|P_n(u)\|_\infty&\leq\|P_n(|u|)\|_\infty\leq\alpha_n(|u|)\leq\frac{\|u\|_1^n}{n!}
\end{aligned}\right.\tag*{(***)}
\]

\uline{Démonstration} On démontre le

\uline{Lemme} Si \inserted{$u$ est opérateur normal} $u=\sum_1^\infty\lambda_i e_i\otimes e'_i$, alors $P_n(u)=\sum\unsure{p_i(u)|\rho_i(u)}e_i\otimes e_i$,

\noindent\begin{tabular}{@{}|p{.92\linewidth}@{}}
où $\unsure{p_i(u)|\rho_i(u)}=\sum_{(I)}'\lambda_{i_1}-\lambda_{i_n}$, $\sum'$ étant étendue \unsure{à|:} tous les \deleted{\unsure{systèmes}} \unsure{\deleted{\unsure{ensembles}}|\inserted{\uline{\unsure{parties}}}}\inserted{$I=\{i_1-i_n\}$} de $n$ indices, \deleted{\unsure{qui tels que $I$} \illegible{2}}\inserted{ne contenant pas $i$}
\end{tabular}

Pour le voir, on calcule successivement (dans la \unsure{base} $\unsure{(e_i)}$) \deleted{\unsure{$u^n$}}$u^n-\alpha_1(u)u^{n-1}$, $u^{\unsure{-|n}}-\alpha_1(u)u^{n-1}+\alpha_2(u)u^{n-1}$ par récurrence on obtient $\sum_{k=0}^{p}u^{n-k}\alpha_k(u)(-1)^k=(-1)^n\sum'\lambda_{i_1}-\lambda_{i_k}$ d'où \deleted{\unsure{$P_n=$ pour tout}} le résultat.

Maintenant 1) devient évident. 2) \unsure{Se|se} vérifie sur les valeurs propres de façon immédiate.

\uline{Remarque} Notons que si $\unsure{\rho_i|p_i}$ sont les valeurs propres de $R(u)$, on a $\unsure{\rho_i|p_i}=\prod_{\substack{\lambda_j\\j\neq i}}(1+\lambda_j)$. d'où $|\unsure{\rho_i|p_i}|\leq\prod_j(1+|\lambda_j|)=\det(1+\unsure{\lvert u\rvert|u})$, on retrouve une

\scanpage{150}
\unsure{partie|partir} de $***$. Mais, si \deleted{$u\to u_{\unsure{0|\infty}}$}\inserted{$1$ n'est pas valeur propre,} on a aussi, \unsure{\deleted{pour \unsure{tout|tt}}|pour \unsure{tout|tt}} \inserted{\unsure{pour la dimension infinie|particulièrement en dimension infinie}} $\unsure{\rho|p}_i\deleted{\unsure{\geq}}=\frac{\det(1+\unsure{\lambda|u})}{1+\lambda_i}\deleted{\unsure{\geq|>}}$ d'où

\[
\boxed{\|R(u)\|_\infty\geq|\det(1+u)|}\quad\text{(op normal)}
\]

(car il existe des $\lambda_i$ arbitrairement petits, donc des $\unsure{\rho|p}_i$ arbitrairement \deleted{\unsure{distants|inférieurs|inférieur} \unsure{par|que}}\inserted{voisins} de $\det(1\unsure{+|-}u)$ !

On \deleted{\illegible{1}} \unsure{prouve} donc que l'inégalité $***$ a) ne peut s'améliorer, car elle devient égalité pour $u\geq0$ !
